\section{The cyclic block-incidence compiler}

\begin{definition}
For an admissible cyclic word \(\gamma=(x_0,\ldots,x_{n-1})\) and an
admissible word \(w=(w_0,\ldots,w_{r-1})\), define
\begin{equation}\label{eq:block-count}
 \Nblk_r(\gamma;w)
 =\#\{0\leq j<n:(x_j,\ldots,x_{j+r-1})=w\},
\end{equation}
where indices are read modulo \(n\).  The vector of all these counts is
denoted by \(\Nblk_r(\gamma)\).
\end{definition}

A width-\(r\) locally constant potential has the form
\(\varphi(x)=v(x_0,\ldots,x_{r-1})\).  Its periodic sum is
\begin{equation}\label{eq:pairing}
 \Birk_\gamma\varphi
 =\sum_{j=0}^{n-1}\varphi(\sigma^j x)
 =\sum_w v(w)\Nblk_r(\gamma;w).
\end{equation}

\begin{lemma}[incidence obstruction]\label{lem:incidence}
If cycles \(\gamma_i\) and integers \(c_i\) satisfy
\begin{equation}\label{eq:incidence-relation}
 \sum_i c_i\Nblk_r(\gamma_i)=0,
\end{equation}
then every width-\(r\) locally constant potential satisfies
\begin{equation}\label{eq:sum-relation}
 \sum_i c_i\Birk_{\gamma_i}\varphi=0.
\end{equation}
Conversely, on any finite list of cycles, a prescribed vector of totals is
realized by a width-\(r\) potential precisely when it annihilates every linear
relation among the corresponding incidence rows.
\end{lemma}

\begin{proof}
The forward statement follows by pairing \eqref{eq:incidence-relation} with
the coefficient vector \(v\).  For the converse, the proposed totals define
a linear functional on the span of the incidence rows exactly when every
row relation is annihilated.  Extend that functional to the ambient finite
word space and read its coordinates as \(v(w)\).
\end{proof}

For widths one and two, the first nontrivial relation is already
\begin{equation}\label{eq:edge-relation}
 \Nblk_r(\gamma_{4a})
 =\Nblk_r(\gamma_1)+\Nblk_r(\gamma_3),
 \qquad r=1,2.
\end{equation}
At width two this simply says that the directed edges of \(\gamma_{4a}\) are
\[
 00,02,21,10,
\]
the disjoint multiset union of the edge \(00\) from \(\gamma_1\) and the
edges \(02,21,10\) from \(\gamma_3\).

The decisive next relation occurs at width three:
\begin{equation}\label{eq:triple-relation}
 \Nblk_3(\gamma_3)+\Nblk_3(\gamma_5)
 =\Nblk_3(\gamma_{4a})+\Nblk_3(\gamma_{4b}).
\end{equation}
Indeed the two sides are the same multiset
\[
 \{002,021,023,100,102,210,231,310\}.
\]
The relation is exact symbolic homology; no numerical orbit data enter it.
